% !TeX program = pdflatex
% UTF-8
%
% coop-writing maximal example, EN-US.
% This file intentionally exercises the public/documented API of v1.5.4.
% It is primarily an EDITING-mode example.
%
% Main modes:
%   editing | submit | publish | acceptingpublish
% Portuguese aliases:
%   edicao  | submeter | publicar | publicaraceitando
%
% Feature/behavior options currently available include:
% comments, anonymize, noanonymize, suggestions, subjects, drafts,
% todos, changecolor, nopdfbookmarks, cwavoidhyperref, toclofttitles.
%
\documentclass[11pt]{article}

\usepackage{iftex}
\ifPDFTeX
  \usepackage[T1]{fontenc}
  \newcommand{\CWEngine}{pdfLaTeX}
\else
  \ifLuaTeX
    \newcommand{\CWEngine}{LuaLaTeX}
  \else
    \newcommand{\CWEngine}{another TeX engine}
  \fi
\fi

\usepackage[english]{babel}
\usepackage[a4paper,margin=28mm]{geometry}
\usepackage{lipsum}
\usepackage[editing]{coop-writing}

% Collaborator families. cwautor is the Portuguese alias of cwnamedef.
\cwnamedef{alice}{blue}{Alice}
\cwautor{bruno}{teal}{Bruno}

% A reusable mode-dependent command for blind review.
\cwblind{\projectname}
  {Open Editorial Project at Example University}
  {Anonymous project}

% Current v1.5.4 defines \rascunho but not \endrascunho.
% This compatibility line can be removed after issue #40 is fixed.
\providecommand{\endrascunho}{\endcwdraft}

% User-facing color configuration available in v1.5.4.
\cwsetanoncolor{violet}
\cwsetsubjectcolor{cyan}
\cwsetmaincolor{yellow}
\cwsetdraftcolor{green}
\cwdefinetodocolor{yellow}

% User-facing replacement strings for anonymous output.
\cwdefanontext{Anonymous institution}
\cwdefanoncitetext{(Anonymous, Year)}
\cwdefanonciteptext{(Anonymous, Year)}
\cwdefanoncitettext{Anonymous (Year)}

% Text used when a suggestion has no explicit comment.
\cwdefattentiontext{Proposed change}

\title{coop-writing maximal example\\EN-US}
\author{Example Author}
\date{\today}

\begin{document}
\maketitle

\section{Version, encoding, and engine}

This source file is UTF-8 and is being compiled with \CWEngine.
The package reports version \coopwritingversion, while
\printcoopwritingversion\ prints its name and version together.

The example contains direct UTF-8 text: Xexéo, naïve, façade, São Paulo,
revisão, informação, ação, and não. These Latin-script examples are part of
the required pdfLaTeX/LuaLaTeX regression coverage proposed in issues \#34
and \#42.

\section{Current localized strings}

The following user-configurable strings are part of the current package surface:

\begin{itemize}
  \item Draft title: \cwdrafttitle.
  \item Comment-list title: \cwcommentstitle.
  \item Subject-list title: \cwsubjtitle.
  \item Citation-needs title: \cwcitationstitle.
  \item Citation request text: \cwpleasecitetext.
  \item Default citation request message: \cwpleasecitemessage.
  \item Citation margin label: \cwpleasecitemarginnote.
\end{itemize}

\section{Basic comments}

The predefined author family creates a comment here
\cwauthor{A comment from the predefined Author collaborator}.
The predefined editor family creates another one
\cweditor[the selected phrase]{A comment from the predefined Editor collaborator}.

The custom Alice family is generated by \texttt{\string\cwnamedef}. First,
a simple comment appears here\alice{Simple Alice comment}. A selected
passage can be annotated as
\alice[this text is selected]{Comment attached to selected text}.

A labeled comment is
\alicer[labeled selected text]{This comment has a stable LaTeX label}{alice-labelled}.
The current reference mechanism produces Comment~\ref{cw:alice-labelled};
issues \#19, \#20, and \#38 discuss a better semantic-ID design.

A struck-through comment is kept for historical context
\alicex{This old comment is visibly resolved}.
The labeled form is also available
\alicerx[old highlighted text]{Resolved labeled comment}{alice-resolved}.

\cwsetcommwarn{!}
This comment has an extra warning mark\alice{The exclamation mark is configured globally}.
\cwsetcommwarn{}

\section{Suggestions, removals, and replacements}

The complete generated family is exercised below.

Alice proposes an insertion: \alicesug{newly proposed text}.

Alice proposes a removal:
\alicerem[This phrase is redundant]{redundant old wording}.

Alice proposes a replacement:
\aliceswap[Improve precision]{an imprecise statement}{a more precise statement}.

The Portuguese synonym generated for every collaborator is also exercised:
\alicetroca[Same operation, Portuguese command name]
{another old formulation}{another new formulation}.

Bruno was defined through \texttt{\string\cwautor}; his generated commands
work in the same way. For example,
\brunosug[Bruno's suggestion]{a suggestion from Bruno}.

\section{Low-level note compatibility API}

The commands in this section have no \texttt{@} in their names and therefore
are technically callable, but normal users should prefer collaborator commands.

\cwnotelabel{Reviewer}
\cwnotemargin{Rv}
A raw note follows\Cwnote{Low-level note created with \string\Cwnote}.

A color-explicit raw note follows
\cwnote[raw-note]{Low-level note created with \string\cwnote}{brown}.

\section{TODO behavior}

The default TODO is intentionally the compact editorial-comment form:
\todo{This TODO must not become an inline box by default.}

The inline form is explicit:
\todo[inline]{This is the intentionally large inline TODO box.}

Issue \#45 records this behavior as a regression requirement.

\section{Citation-needed annotations}

This factual claim needs support.\pleasecite

This second claim asks for a specific kind of source.
\pleasecite[Prefer a recent peer-reviewed systematic review.]

The Portuguese alias works too.\favorcitar[The alias should create the same semantic item.]

\section{Draft material}

\begin{cwdraft}[Alternative argument]
This block is provisional. It is visible in editing mode and should disappear
from clean modes.
\end{cwdraft}

The Portuguese environment alias is demonstrated with the temporary
v1.5.4 end-macro compatibility line from the preamble:

\begin{rascunho}[Portuguese alias]
This is the same draft mechanism reached through the Portuguese alias.
\end{rascunho}

\section{Anonymization}

The institution is wrapped directly as
\cwanon{Example University}.

A reusable blind-review macro expands here as \projectname.

The generic citation wrapper gives \cwanoncite{demo-ref}.
The parenthetical-compatible wrapper gives \cwanoncitep[p.~42]{demo-ref}.
The textual-compatible wrapper gives \cwanoncitet{demo-ref}.

In editing mode these items are shown in the configured anonymization color;
in submit mode they use the replacement strings configured in the preamble.

\section{Subjects and paragraph main ideas}

\cwsubject{First subject}
This paragraph begins with a subject marker.

\cwassunto{Subject through the Portuguese alias}
This paragraph demonstrates \texttt{\string\cwassunto}.

\cwmain{This sentence states the main idea of the paragraph.}
The remainder of the paragraph explains and supports it.

The emphasis macro can be used directly:
\cwmainemphasis{direct use of the current main-emphasis renderer}.

It can also be reconfigured with ordinary LaTeX:
\renewcommand{\cwmainemphasis}[1]{\textbf{#1}}
\cwmain{This main idea is rendered in bold after redefining the emphasis hook.}

\section{Mode-dependent files}

The three-way input selector follows. In this editing build, the first file is used:

\cwinput
  {example-fragments/editing-input}
  {example-fragments/submit-input}
  {example-fragments/publish-input}

The mode-specific input commands are all present in the source:

\cwinputediting{example-fragments/editing-input}
\cwinputsubmit{example-fragments/submit-input}
\cwinputpublish{example-fragments/publish-input}

The three-way include selector follows:

\cwinclude
  {example-fragments/editing-include}
  {example-fragments/submit-include}
  {example-fragments/publish-include}

The mode-specific include commands are all present in the source. Only the
editing one executes in this build:

\cwincludeediting{example-fragments/editing-input}
\cwincludesubmit{example-fragments/submit-input}
\cwincludepublish{example-fragments/publish-input}

\section{Text-color helpers}

Normal color. \cwcolor{blue}This text is changed with
\texttt{\string\cwcolor}.\cwrestorecolor\ Normal color again.

The lower-level save helper can be used separately:
\cwsavecolor
\color{magenta}magenta after an explicit \texttt{\string\color};
\cwrestorecolor
the saved color is restored.

\section{A little ordinary document content}

The package is intended to coexist with ordinary LaTeX structures, equations,
floats, citations, and long prose. For example,
\[
  f(x)=x^2+1
\]
is ordinary mathematics and should not be affected by the editorial layer.

\lipsum[2]

\begin{thebibliography}{1}
\bibitem{demo-ref}
A. Author.
\newblock A demonstration reference.
\newblock \emph{Journal of Examples}, 1(1):1--10, 2026.
\end{thebibliography}

\clearpage
\section*{Editorial lists: canonical commands}

\listofcomments
\listofcitationneeds
\listofsubjects

\clearpage
\section*{Editorial lists: compatibility aliases}

The following intentionally duplicate the lists above so that the maximal
example exercises the legacy/underlying public aliases as well.

\listofcomentario
\listofcomentarioref
\listofsubject
\listofassunto

\end{document}
